% 附录：AI 工具使用情况说明（国赛 / 中文竞赛通用四问式模板）
% 说明：请全部填写真实记录（取自 work/iteration_log.md 与审稿修改记录），严禁编造。

\subsection*{一、AI 工具使用清单与分工}
本参赛队在遵循竞赛规则的前提下，仅将 AI 工具作为辅助效率工具用于文献初筛、代码框架编写、图表脚本调试与语言校对，核心建模假设、数学推导及最终结论均由参赛队员独立完成并逐项审定。具体分工如下表所示：

\begin{table}[H]
\centering
\caption{AI 工具辅助环节与人工审核分工表}
\begin{tabular}{p{3cm}p{4cm}p{4cm}p{3.5cm}}
\toprule
工作环节 & 使用工具/版本 & AI 辅助内容 & 人工核验与决策 \\
\midrule
题面解析与校验 & [模型名称] & 双通道核对题面条件与表格 & 人工逐页对照原题 PDF 确认 \\
模型编程实现 & [模型名称] & 辅助编写求解器与绘图脚本 & 独立约束检查器逐条复算 \\
论文排版与润色 & [模型名称] & LaTeX 语法检查与表达梳理 & 人工通读并核对全部数字 \\
\bottomrule
\end{tabular}
\end{table}

\subsection*{二、输出核验机制}
为防止 AI 产生数学错误、代码幻觉或引用失真，本队建立了三层机械核验流程：
\begin{enumerate}
  \item \textbf{形式化与约束复核}：先人工审定数学形式化文档，再编写独立于求解器逻辑的约束检查脚本，对全部约束逐条计算违约量。
  \item \textbf{数值唯一溯源}：正文、摘要与图表中的所有实验数字统一由求解产物汇编并运行双向一致性校验脚本，杜绝手抄或估算数字。
  \item \textbf{文献与版面校验}：通过 Crossref/arXiv 核验全部参考文献元数据，并对插图与成稿 PDF 进行像素级与字号下限审计。
\end{enumerate}

\subsection*{三、人工纠正 AI 失误的具体实例}
在建模与写作过程中，本队发现并纠正了以下真实问题（详见迭代日志）：
\begin{itemize}
  \item \textbf{实例 1（建模/约束纠正）}：[填写具体发生的公式方向、边界条件或单位换算纠偏]
  \item \textbf{实例 2（数据/代码纠正）}：[填写具体发现的索引错位、不可行解或特征泄露修复]
  \item \textbf{实例 3（写作/数字纠正）}：[填写具体拦截的数字不一致、表注错误或排版溢出]
\end{itemize}

\subsection*{四、责任声明}
本参赛队声明：上述 AI 使用情况完全属实；论文中的所有数学模型、程序代码、实验数据、参考文献及最终结论均已通过参赛队员的人工复核与验证，本队对提交作品的学术规范性与全部内容负最终全责。
